\section{Which experiment distinguishes a sign?}
Two classical universes can have the same spatial metric and opposite homogeneous axial torsion. A spinor measures both the metric and the connection. Tuning a signed torsion until its spatial Dirac operator has a zero mode is one experiment. Measuring the occupation of a prepared state after a prescribed expansion is another. These experiments require distinct inputs even when the charge has the same origin. We first fix the sample space for each question, then compute examples in which the answer changes with the spin lift, aspect ratio or preparation.

The auxiliary connection of a dynamical Nieh--Yan coupling leads to a massless pseudoscalar in the bosonic action, a mechanism analyzed by Calcagni and Mercuri \cite{calcagni}. Paper~04 \cite{p04} gives uncharged Dirac data on selected compact geometries, and Paper~07 \cite{p07} fixes the Lorentzian coframe and boundary conventions for constant Nieh--Yan coupling. Here a spatially homogeneous dynamical scalar supplies the conserved charge. We insert an external, minimally coupled spinor to define a probe. A geometric zero mode, a regularized vacuum-energy difference and an occupation probability are different outputs of this choice.

The principal mathematical question has an exact answer on the reference Sol family. A hyperbolic integral monodromy \(A\) assigns a nonzero integer form value \(F(v)\) to each nonzero Fourier vector. Section~\ref{sec:sol} proves that equality of the least positive value and the least absolute negative value of \(F\) is equivalent to equality of the first nonzero-fibre charge magnitudes. The proof works for every positive hyperbolic trace. The additional agreement with topological chirality is an enumeration for traces at most 60. On a chiral Sol example with \(\eta=0\), the magnitude at which negative charge first produces a kernel differs from the positive one by more than seven charge units.

Ammann and B\"ar \cite{ammann}, Gornet and Richardson \cite{gornet}, and Pf\"affle \cite{pfaffle} computed unshifted spectra and eta invariants on Heisenberg or flat spaces. B\"ar \cite{baer} and Boldt and Lauret \cite{boldt} provide representation-theoretic tools for spherical quotients. Our comparisons use these spectra with a signed charge shift, keeping track of which lift and which subspace the experiment retains. On Sol, Lin's unshifted kernel theorem \cite[Lemma 4.2]{lin}, Marcolli's adiabatic-limit construction \cite{marcolli}, and the arithmetic Laplace multiplicities of Bolsinov, Dullin and Veselov \cite{bolsinov} concern related operations. Here the datum is the least positive and negative spin level at fixed metric on the nonzero-fibre space.

Section~\ref{sec:device} replaces a family of static operators by actual evolution. The two signs have the same metric, laboratory projectors and elapsed interval. A real Dirac mass couples the two Weyl components. This makes the marked shell sensitive to a torsion term which otherwise contributes only a scalar phase to massless two-band transport.
